
\chapter{Chapter 1. Introduction}

\section{Motivation}

Dense vector retrieval has become a core dependency of modern AI systems. Every
retrieval-augmented generation (RAG) pipeline~\cite{lewis2020retrieval}, enterprise search
engine, and question-answering service relies on the ability to find semantically relevant
content in large corpora at query time. Unlike sparse methods such as TF-IDF or BM25,
dense retrieval encodes queries and passages as continuous vector embeddings using bi-encoder
architectures~\cite{reimers2019sbert, karpukhin2020dense}, and measures semantic similarity
by cosine or dot-product distance. This enables retrieval that generalizes well beyond exact
keyword matching, which has driven rapid adoption across search, recommendation, and
generative AI applications.

As retrieval corpora scale to tens of millions or billions of passages, the memory footprint
of stored embeddings becomes a critical infrastructure bottleneck. A standard model such as
BGE-base-en-v1.5~\cite{xiao2023c} produces 768-dimensional float32 vectors, consuming
3,072~bytes per passage. Indexing one hundred million passages therefore requires
approximately 286~GB; at one billion passages this grows to approximately 2.9~TB.

Dense retrieval systems built on libraries such as FAISS~\cite{johnson2021billion} typically
require the embedding index to reside in random-access memory (RAM) for interactive query
latencies: DRAM access is on the order of 100~ns, roughly three orders of magnitude faster
than NVMe SSD random reads. Serving these in-memory indices requires memory-optimized cloud
instances; the AWS x2gd family, for example, costs approximately \$3.80 per GB per
month~\cite{aws2024pricing}. Table~\ref{tab:index-cost-intro} shows how this cost scales
with embedding dimension and corpus size across widely-used models.

\begin{table}[ht]
\caption[Float32 index size and monthly instance cost by embedding dimension]{Float32
embedding index size and estimated monthly cost on memory-optimized AWS x2gd instances
(\$3.80/GB/month) at two corpus scales~\cite{shakir2024embedding, aws2024pricing}.
Sizes are reported in GB (SI).}
\label{tab:index-cost-intro}
\begin{center}
\resizebox{\columnwidth}{!}{
\begin{tabular}{| c | c | c | c |}
\hline
\textbf{Dim} & \textbf{Example Models} & \textbf{100M Passages} & \textbf{1B Passages} \\ [0.5ex]
\hline
384  & all-MiniLM-L6-v2, bge-small-en-v1.5          & 143 GB / \$543/mo    & 1,431 GB / \$5,435/mo  \\
\hline
768  & bge-base-en-v1.5, nomic-embed-text-v1         & 286 GB / \$1,087/mo  & 2,861 GB / \$10,871/mo \\
\hline
1024 & bge-large-en-v1.5, mxbai-embed-large-v1       & 381 GB / \$1,449/mo  & 3,815 GB / \$14,495/mo \\
\hline
1536 & text-embedding-3-small                        & 572 GB / \$2,174/mo  & 5,722 GB / \$21,743/mo \\
\hline
3072 & text-embedding-3-large                        & 1,144 GB / \$4,348/mo & 11,444 GB / \$43,487/mo \\
\hline
\end{tabular}
}
\end{center}
\end{table}

As Table~\ref{tab:index-cost-intro} shows, a billion-passage index at 768 dimensions already
exceeds \$10,800 per month, and the cost more than quadruples when moving to 3,072-dimensional
models. These figures represent the minimum viable infrastructure: the memory requirement
dictates the instance size, and therefore the total cost, including compute. The problem is
growing worse as the field trends toward higher-dimensional embeddings (models such as
text-embedding-3-large at 3,072 dimensions are increasingly common) and as RAG-based
applications push corpus sizes into the billions of passages.

Beyond cost, larger embeddings also hurt speed: every query must read each document vector
from memory, so the bigger the vectors, the fewer queries the system can serve per second. For resource-constrained
environments such as on-device or edge deployments, multi-terabyte indices are simply
infeasible regardless of budget. Embedding compression is therefore not merely a storage
optimization; it is a prerequisite for deploying dense retrieval at scale.

Even simple compression approaches already show real promise: replacing each 32-bit float
with a single bit (binary quantization) shrinks the index by $32\times$ and retains much of
the retrieval quality when applied to a strong pre-trained model~\cite{shakir2024embedding}.
Compression can be pushed further by stacking multiple techniques. The open questions are how
far compression can go before retrieval quality degrades substantially, and whether
incorporating compression awareness into fine-tuning produces better representations than
applying compression post-hoc. This thesis investigates both.

\section{Embedding Compression Landscape}

This thesis organizes embedding compression methods along two orthogonal axes.
The first axis is \textit{when} compression is applied: \textbf{post-hoc} methods
are applied to the generated embeddings without modifying the model, while \textbf{training-time}
methods incorporate compression awareness into the model's fine-tuning objective.
The second axis is \textit{how} memory is reduced: \textbf{precision reduction} decreases
the number of bits per dimension, \textbf{dimensionality reduction} decreases the number
of dimensions kept, and \textbf{learning to hash} fits a codebook that maps embeddings to discrete codes directly,
rather than applying scalar quantization per dimension.
Table~\ref{tab:compression-taxonomy} places the methods studied in this thesis along
these two axes.

\begin{table}[ht]
\caption[Organization of embedding compression methods studied in this thesis]{Organization of embedding compression methods studied in this thesis, by when compression is applied and how memory is reduced.
Abbreviations: INT8 = 8-bit integer quantization;
PQ = Product Quantization;
BAT = Binary-Aware Training (STE and Annealed Tanh are two training variants);
STE = Straight-Through Estimator; MRL = Matryoshka Representation Learning.}
\label{tab:compression-taxonomy}
\begin{center}
\begin{tabular}{| l | c | c | c |}
\hline
 & \textbf{Precision Reduction} & \textbf{Dim.\ Reduction} & \textbf{Learning to Hash} \\ [0.5ex]
\hline
\textbf{Post-hoc}       & INT8, Binary, TurboQuant & ---  & PQ            \\
\hline
\textbf{Training-time}  & BAT (STE / Annealed Tanh)  & MRL  & ---  \\
\hline
\end{tabular}
\end{center}
\end{table}

Each method reduces embedding memory through a distinct mechanism.
\textbf{8-bit integer (INT8)} quantization maps each float32 value to an 8-bit integer,
yielding $4\times$ compression with minimal quality loss.
\textbf{Binary quantization} maps each dimension to a single bit ($32\times$ compression),
enabling fast Hamming-distance similarity search.
\textbf{Product Quantization (PQ)}~\cite{jegou2010product} partitions each embedding into
fixed-width sub-vectors and quantizes each to a small learned codebook, achieving high
compression ratios with better quality retention than scalar binary quantization.
\textbf{TurboQuant}~\cite{zandieh2025turboquant} applies a random rotation to each embedding before scalar quantization, distributing distortion uniformly across dimensions to approach information-theoretic compression bounds.
\textbf{Binary-Aware Training (BAT)} fine-tunes the model with an inserted sign function
to produce binary-friendly representations; two gradient approximation variants are studied:
the Straight-Through Estimator (STE)~\cite{bengio2013estimating}, which passes gradients
unchanged through the discontinuity, and Annealed Tanh, which replaces the STE with a
temperature-annealed $\tanh(\beta \cdot x)$ surrogate that smoothly transitions toward the
sign function as training progresses.
\textbf{Matryoshka Representation Learning (MRL)}~\cite{kusupati2022matryoshka} trains
nested subspace representations so that any leading $k$-dimensional prefix suffices for
retrieval, enabling post-hoc truncation to any desired dimensionality at inference time.

The fundamental tension in embedding compression is between compression depth and retrieval
quality. Every method in Table~\ref{tab:compression-taxonomy} trades information loss for
a smaller memory footprint: reducing bits per dimension discards numerical precision,
reducing dimensions discards representational capacity, and learned codes introduce
quantization error. At moderate compression (INT8 at $4\times$), quality losses are
negligible. At aggressive compression (binary embeddings at $32\times$, heavily truncated
MRL representations, or stacked combinations), degradation becomes more substantial, and
the shape of this trade-off varies significantly across methods and compression levels.

\section{Problem Statement}

The most fundamental open question in embedding compression is whether training-time methods
actually outperform post-hoc compression, a comparison that has never been made under
controlled conditions. BPR~\cite{yamada2021efficient} and QAMA~\cite{qama2025} both
demonstrate that training-time binarization works well in their respective settings, while
post-hoc studies~\cite{shakir2024embedding} also report strong results; yet these bodies
of work use different base models, different training corpora, and different benchmarks.
When everything changes at once, a method that appears to win may simply reflect a better
starting point rather than a genuinely better compression strategy. CoRECT~\cite{correct2025},
the most recent large-scale evaluation, takes a step toward systematic comparison, but its
primary focus is how compression holds as corpus size scales, it covers only a limited set
of post-hoc techniques, and it leaves training-time methods entirely out of scope. A
secondary open question concerns stacking: while post-hoc combinations such as MRL with
binary quantization have received some attention~\cite{correct2025}, broader combinations
involving training-time methods (such as MRL jointly trained with BAT or PQ) remain
unstudied.

This thesis addresses that gap through a controlled empirical study. We fine-tune three
encoder architectures (BGE-base-en-v1.5~\cite{xiao2023c}, RoBERTa-base~\cite{liu2019roberta},
and MPNet-base~\cite{song2020mpnet}) on the same multi-source retrieval corpus of
4.3~million query-passage pairs, and evaluate all compression methods on the same benchmark.
By holding training data and evaluation fixed and varying only the compression method, any
difference in retrieval quality can be attributed to the compression strategy itself rather
than to the choice of model. Post-hoc and training-time methods are evaluated side by side,
and MRL is systematically combined with five compression techniques to characterize the
stacked compression frontier. All results are reported as Pareto curves on the NanoBEIR
benchmark~\cite{thakur2021beir}, directly answering the practitioner's question: given a
memory budget, which approach should you use?

\section{Research Questions}

This thesis is organized around four research questions:

\begin{enumerate}

    \item \textbf{Training-time versus post-hoc compression.}
    Does embedding quantization-aware training (such as BAT with STE or Annealed Tanh,
    or MRL) yield better retrieval quality than post-hoc compression applied to a
    float-trained baseline at equivalent embedding sizes?

    \item \textbf{Binarization training strategy.}
    Between STE and Annealed Tanh Binarization, which gradient surrogate produces
    better binary embeddings, and how sensitive are results to the annealing rate?

    \item \textbf{Compression stacking with MRL.}
    Does combining MRL with other compression methods (binary quantization, PQ,
    TurboQuant, and BAT) yield better Pareto efficiency than either technique
    applied independently?

    \item \textbf{Cross-architecture generalizability.}
    Do the above findings hold consistently across different encoder architectures
    (BGE-base, RoBERTa-base, and MPNet-base), or are they specific to a particular model?

\end{enumerate}

\section{Contributions}

This thesis makes the following contributions to the field of efficient dense text retrieval:

\begin{enumerate}

    \item \textbf{Systematic cross-method comparison.}
    Post-hoc and training-time compression methods are evaluated side by side under
    identical training data and benchmark, across three encoder architectures (BGE-base,
    RoBERTa-base, MPNet-base). Both single methods and stacked combinations are covered,
    spanning binary, scalar, product quantization, and dimensionality reduction at multiple
    operating points. Prior work evaluates each method under different setups; this study
    provides a direct, fair comparison.

    \item \textbf{Binarization strategy comparison.}
    STE and Annealed Tanh have each been used for retrieval binarization in prior
    work~\cite{bengio2013estimating, yamada2021efficient}, but never compared directly.
    Both are trained under identical conditions to isolate the effect of the gradient
    surrogate on retrieval quality. A sweep of the temperature schedule parameter measures
    how sensitive Annealed Tanh is to the annealing rate.

    \item \textbf{Compression stacking with MRL.}
    MRL is combined with five compression techniques: post-hoc binary quantization, PQ,
    TurboQuant, BAT (STE), and Annealed Tanh. Results establish which combinations achieve
    the best Pareto trade-offs and whether jointly training MRL with binarization
    outperforms applying the two techniques independently.

    \item \textbf{Reproducible experimental pipeline.}
    A DDP fine-tuning setup over seven heterogeneous retrieval datasets with
    domain-weighted batch sampling provides the controlled training baseline. An embedding
    caching and transform pipeline decouples encoding from evaluation, enabling any
    compression variant to be benchmarked without re-encoding. Results are surfaced through
    an interactive Plotly dashboard with pairwise Wilcoxon signed-rank and paired t-test
    significance testing across all model pairs.

\end{enumerate}

\section{Thesis Organization}

The remainder of this thesis is organized as follows. Chapter 2 provides background on dense
retrieval, contrastive learning objectives, and the embedding compression methods compared in
this work. Chapter 3 describes the training and evaluation datasets, including the multi-source
training corpus construction and the NanoBEIR benchmark. Chapter 4 details the training setup
and the fine-tuned float32 baseline that serves as the anchor for all downstream comparisons.
Chapter 5 presents the compression methods, both post-hoc and training-time, including the
derivation of the Annealed Tanh Binarization approach and its relationship to the
STE-based BAT baseline. Chapter 6 reports the experimental results, including NanoBEIR
evaluations and Pareto analysis across all methods.
Chapter 7 concludes with a summary of findings, practical guidance on method selection at
given memory budgets, and directions for future work.
